% Section 9.1, from lectures/L09/appendix/a_flowcharts_and_branches.md.
%
% Typesetting only: in the listings of 9.1.5 and 9.1.6 the comments start four columns further left
% than in the lecture, so that the longest line fits the listing's width.

\section{Flödesplaner och villkorliga hopp}
\label{c9:app:a}

\subsection{Flödesplaner}
\label{c9:a1}
Programmen hittills har gått rakt uppifrån och ned, en instruktion efter en annan. Riktiga program
måste \textbf{välja}: tänd värmen om det är kallt, annars inte. Innan ett sådant program skrivs är
det värt att rita det, som en \textbf{flödesplan}.

En flödesplan beskriver vad programmet ska göra, i vilken ordning, och var det väljer väg, utan att
bry sig om vilka instruktioner som behövs. Den ritas med fem symboler:

\bookfigure{l09_symbols}{De fem symbolerna i en flödesplan.}{c9:fig:symbols}

\begin{booktable}{@{}p{6.4em}Lp{8.6em}@{}}
  \thead{Symbol} & \thead{Betyder} & \thead{Exempel} \\
  \midrule
  Start/stopp & Var programmet börjar och slutar & Start \\
  Process & Något programmet gör & \code{summa = summa + 1} \\
  Beslut & En fråga med två utgångar, Ja och Nej & \code{r16 = 0?} \\
  In/utmatning & Läsa en ingång eller skriva en utgång & \code{PORTB = mönster} \\
  Subrutin & Ett delprogram som anropas; tas upp i kapitel~\ref{c10:ch} & \code{delay} \\
\end{booktable}

\noindent Fyra regler gör en flödesplan lätt att läsa, och lätt att översätta till ett program:
\begin{itemize}
  \item \textbf{En start}, och flödet går uppifrån och ned. Pilar som går uppåt är loopar.
  \item \textbf{Varje beslut har exakt två utgångar}, märkta Ja och Nej.
  \item \textbf{Skriv vad, inte hur.} \code{Värme på} är bättre än \code{sbi PORTB, 0}; vilken
    instruktion det blir bestäms när programmet skrivs.
  \item \textbf{Grenar som delar sig ska mötas igen.} Mer om det i \secref{c9:a8}.
\end{itemize}

\subsection{Från flödesplan till program}
\label{c9:a2}
När flödesplanen är klar översätts den ruta för ruta:
\begin{itemize}
  \item en \textbf{process} eller en \textbf{in/utmatning} blir en eller några få instruktioner;
  \item ett \textbf{beslut} blir en jämförelse följd av ett villkorligt hopp (\secref{c9:a3}
    och~\ref{c9:a4});
  \item varje ställe dit en \textbf{pil pekar} från mer än ett håll får en \textbf{etikett}, så att
    ett hopp kan hamna där.
\end{itemize}

Det sista är nyckeln. Ett program kör uppifrån och ned tills en hoppinstruktion skickar det någon
annanstans, och dit måste det finnas en etikett att hoppa till.

\subsection{Jämförelse: \code{cp}, \code{cpi} och \code{tst}}
\label{c9:a3}
Ett beslut i en flödesplan är en fråga om två tal: är de lika, är det ena mindre? På AVR ställs
frågan med en \textbf{jämförelse}:

\begin{booktable}{@{}p{5.6em}p{7em}L@{}}
  \thead{Instruktion} & \thead{Exempel} & \thead{Gör} \\
  \midrule
  \code{cp} & \code{cp r16, r17} & Beräknar \code{r16 - r17} och sparar bara flaggorna. \\
  \code{cpi} & \code{cpi r16, 20} & Beräknar \code{r16 - 20} och sparar bara flaggorna. Bara
    \code{r16}--\code{r31}. \\
  \code{tst} & \code{tst r16} & Sätter \code{Z} om \code{r16} är noll och \code{N} om bit~7 är
    satt. \\
\end{booktable}

\noindent \textbf{En jämförelse är en subtraktion vars resultat kastas bort.} \code{cp r16, r17}
räknar ut \code{r16 - r17} precis som \code{sub} gör, men lägger inte svaret någonstans: \code{r16}
är oförändrat efteråt. Det enda som blir kvar är flaggorna, och de säger allt man behöver veta om de
två talen:

\begin{booktable}{@{}p{12em}p{10em}L@{}}
  \thead{Om \code{r16} och \code{r17} är} & \thead{blir \code{r16 - r17}} &
    \thead{och flaggorna} \\
  \midrule
  lika & noll & \code{Z = 1}, \code{C = 0} \\
  \code{r16} större, utan tecken & positivt, inget lån & \code{Z = 0}, \code{C = 0} \\
  \code{r16} mindre, utan tecken & ett lån behövs & \code{Z = 0}, \code{C = 1} \\
\end{booktable}

\noindent Flaggorna och deras betydelse kommer från \secref{c8:a6}. Det är det flaggorna är till
för.

\textbf{Är du osäker på vilken väg en jämförelse läses}, skriv ut subtraktionen. \code{cp r16, r17}
är \code{r16 - r17}, alltid i den ordningen.

\subsection{Villkorliga hopp}
\label{c9:a4}
Ett \textbf{villkorligt hopp} hoppar till en etikett om en flagga har ett visst värde, och
fortsätter annars med nästa instruktion. Tillsammans med en jämförelse blir det ett beslut. De
vanligaste:

\begin{booktable}{@{}p{3.6em}Lp{4.6em}@{}}
  \thead{Hopp} & \thead{Efter \code{cp a, b}: hoppar om} & \thead{Läser} \\
  \midrule
  \code{breq} & $a = b$ (\emph{branch if equal}) & \code{Z = 1} \\
  \code{brne} & $a \neq b$ (\emph{branch if not equal}) & \code{Z = 0} \\
  \code{brlo} & $a < b$, utan tecken (\emph{lower}) & \code{C = 1} \\
  \code{brsh} & $a \geq b$, utan tecken (\emph{same or higher}) & \code{C = 0} \\
  \code{brlt} & $a < b$, med tecken (\emph{less than}) & \code{S = 1} \\
  \code{brge} & $a \geq b$, med tecken (\emph{greater or equal}) & \code{S = 0} \\
  \code{brmi} & resultatet negativt (\emph{minus}) & \code{N = 1} \\
  \code{brpl} & resultatet positivt eller noll (\emph{plus}) & \code{N = 0} \\
\end{booktable}

\noindent Hela listan finns på referensbladet i bilaga~\ref{app:avr}. Två saker till är värda att
veta:
\begin{itemize}
  \item \textbf{\code{rjmp}} är det \textbf{ovillkorliga} hoppet: det hoppar alltid.
  \item \textbf{Ett villkorligt hopp tar 1 klockcykel om det inte hoppar och 2 om det hoppar},
    eftersom processorn då måste hämta en annan instruktion än den som stod på tur. Det spelar
    roll för tidsfördröjningarna i \secref{c9:app:b}.
\end{itemize}

Det finns inga hopp för \enquote{större än} och \enquote{mindre än eller lika med}. De behövs inte:
$a > b$ är detsamma som $b < a$, så man byter plats på operanderna i \code{cp}, eller jämför med ett
tal ett steg högre: \code{r16 > 24} är \code{r16 ≥ 25}.

\subsection{Om, annars}
\label{c9:a5}
Det vanligaste beslutet är \textbf{om-annars} (\emph{if-else}): om villkoret gäller, gör en sak,
annars en annan. Här är värmaren i \code{compare_branch.asm}, först som flödesplan:

\bookfigure{l09_if_else}{Flödesplan för värmaren: ett om-annars.}{c9:fig:if_else}

\needspace{10\baselineskip}
\noindent och sedan som program:

\begin{asmcode}
    cpi r16, LIMIT              ; Compare: computes r16 - 20, keeps only the flags.
    brlo heater_on              ; Branch if lower (C = 1): temperature < 20.
    cbi PORTB, 0                ; Otherwise: heater off,
    rjmp done                   ; and jump past the other branch.
heater_on:
    sbi PORTB, 0                ; Heater on.
done:
    rjmp done
\end{asmcode}

\noindent Följ de två vägarna:
\begin{itemize}
  \item \textbf{Temperaturen är 18.} \code{cpi} räknar $18 - 20$, som kräver ett lån, så
    \code{C = 1}. \code{brlo} hoppar till \code{heater_on}, och värmen slås på.
  \item \textbf{Temperaturen är 22.} $22 - 20$ kräver inget lån, \code{C = 0}. \code{brlo} hoppar
    inte, och programmet fortsätter med \code{cbi}, som slår av värmen. Sedan hoppar
    \code{rjmp done} \textbf{förbi} \code{heater_on}.
\end{itemize}

\textbf{Det sista \code{rjmp} är lätt att glömma, och då blir programmet fel.} Utan det skulle
programmet efter \code{cbi} fortsätta rakt ned in i \code{heater_on} och slå på värmen igen. I
flödesplanen ser det ut som att grenarna är åtskilda; i programmet ligger de efter varandra, och
det är hoppen som håller isär dem.

\subsubsection{Bara om}
Utan annars-gren blir det enklare: hoppa \textbf{förbi} det som bara ska göras ibland. Då hoppar
man på det \textbf{motsatta} villkoret:

\begin{asmcode}
    cpi r16, 30
    brlo not_too_hot                ; Below 30: skip the alarm.
    sbi PORTB, 7                    ; 30 or above: sound the alarm.
not_too_hot:
\end{asmcode}

\noindent Flödesplanen säger \enquote{om r16 $\geq$ 30, larma}, och programmet säger
\enquote{om r16 $<$ 30, hoppa över larmet}. Det är samma sak, och det är det vanligaste
mönstret i assembler.

\subsection{Utan tecken eller med tecken}
\label{c9:a6}
Samma jämförelse kan ge olika svar beroende på om talen läses med eller utan tecken
(\secref{c8:a5}). Jämför $-5$ med 3:

\begin{asmcode}
    ldi r16, -5                 ; 0xFB: 251 without sign, -5 with.
    ldi r17, 3
    cp r16, r17                 ; Unsigned 251 > 3: C = 0. Signed -5 < 3: S = 1.
\end{asmcode}

\noindent Efter \code{cp} är \code{C = 0}, så \code{brlo} hoppar \textbf{inte}: utan tecken är 251
inte mindre än 3. Men \code{S = 1}, så \code{brlt} hoppar: med tecken är $-5$ mindre än 3. Välj
hopp efter vad talen betyder:
\begin{itemize}
  \item \textbf{utan tecken}, räknare, antal, portvärden: \code{brlo} och \code{brsh};
  \item \textbf{med tecken}, tal som kan bli negativa, som en utomhustemperatur: \code{brlt} och
    \code{brge}.
\end{itemize}

Värmaren i \secref{c9:a5} använder \code{brlo}. Det fungerar så länge temperaturen aldrig är
negativ, vilket gäller inomhus. Vid $-5$ grader läser \code{brlo} talet som 251, bedömer att det är
varmt nog, och slår av värmen. En termostat för minusgrader behöver \code{brlt}; se
övning~\ref{c9:ex:12}.

\subsection{Flerval}
\label{c9:a7}
Ibland finns fler än två alternativ: ett läge 0, 1, 2 eller 3, som var och ett ska ge ett eget
mönster på lysdioderna. Det byggs som en \textbf{kedja av beslut}, där varje beslut testar ett
alternativ:

\bookfigure{l09_menu}{Flödesplan för ett flerval med en gemensam utgång.}{c9:fig:menu}

Programmet \code{menu.asm} följer flödesplanen rad för rad:

\begin{asmcode}
    cpi r16, 0
    breq mode_0
    cpi r16, 1
    breq mode_1
    cpi r16, 2
    breq mode_2
    cpi r16, 3
    breq mode_3
    ldi r17, 0b01010101             ; None of them: the error pattern.
    rjmp show
mode_0:
    ldi r17, 0b00000000
    rjmp show
    ...
mode_3:
    ldi r17, 0b11110000
show:
    out PORTB, r17                  ; The one place every alternative ends up.
\end{asmcode}

\noindent Varje alternativ gör bara en sak, lägger sitt mönster i \code{r17}, och hoppar sedan till
samma ställe, \code{show}. Det sista alternativet behöver inget \code{rjmp show}, eftersom
\code{show} kommer direkt efter.

Tänk på \textbf{alternativet som inte finns med}: vad ska hända om läget är 7? En flödesplan utan
svar på den frågan ger ett program som gör något ingen har bestämt. Här visar det ett felmönster.

\subsection{Att arbeta om en flödesplan}
\label{c9:a8}
Den första flödesplanen man ritar blir sällan den bästa. Ett vanligt första försök för flervalet
ser ut så här: varje ja-gren får en egen \code{PORTB = mönster} och en egen stopp-ruta. Det
fungerar, men flödesplanen får fyra slut, och programmet får fyra kopior av \code{out} och fyra
slutloopar. Den dag något ska ändras i slutet, en fördröjning, en loop tillbaka till början, måste
det göras på fyra ställen, och ett av dem glöms.

Att \textbf{arbeta om} flödesplanen betyder att flytta det som är gemensamt till ett ställe:
\begin{enumerate}
  \item \textbf{Leta efter rutor som gör samma sak} i flera grenar. Här är det
    \code{PORTB = mönster} och stopp.
  \item \textbf{Låt grenarna bara göra det som skiljer dem åt}, att välja mönster, och låt dem
    mötas.
  \item \textbf{Rita den gemensamma delen en gång}, efter mötespunkten.
\end{enumerate}

Resultatet är flödesplanen i \secref{c9:a7}: en väg in, en väg ut, och varje alternativ en enda
ruta. Det är den form som blir ett kort och läsbart program, och den form som går att bygga vidare
på. I Labb~3 \hyperref[lab3:d17]{del~17} gör du samma omarbetning själv.

\subsection{Vanliga fel}
\label{c9:a9}
\begin{itemize}
  \item \textbf{Flaggorna skrivs över mellan jämförelsen och hoppet.} Nästan alla
    räkneinstruktioner ändrar flaggorna. Lägg hoppet direkt efter \code{cp} eller \code{cpi}.
    \code{ldi}, \code{mov} och \code{out} ändrar inga flaggor och är ofarliga emellan;
    \code{inc}, \code{dec} och \code{add} är det inte.
  \item \textbf{Fel sorts hopp.} \code{brlo} för tal med tecken, eller \code{brlt} för tal utan,
    fungerar för de flesta värden och blir fel för några. Se \secref{c9:a6}.
  \item \textbf{Det saknade \code{rjmp}} efter den första grenen i ett om-annars, så att programmet
    ramlar in i den andra grenen. Se \secref{c9:a5}.
  \item \textbf{Omvänd jämförelse.} \code{cp r16, r17} följt av \code{brlo} hoppar om
    \code{r16 < r17}, inte tvärtom. Skriv ut subtraktionen om du är osäker.
  \item \textbf{Hoppet når inte fram.} Ett villkorligt hopp når bara ungefär 64 instruktioner
    bort. Längre bort ger ett felmeddelande när programmet byggs, och löses genom att hoppa på det
    motsatta villkoret förbi ett \code{rjmp}, som når ungefär 2\,000 instruktioner åt vardera
    hållet.
\end{itemize}
